% ============================================================
% Section 58: 弗仑克尔缺陷的平衡浓度
% ============================================================
\newpage
\section{固体理论：弗仑克尔缺陷的平衡浓度}
\label{sec:frenkel-defect}

\begin{modelbox}[题目：弗仑克尔缺陷数目与温度]
(1) 证明热平衡时，晶体中弗仑克尔缺陷的数目 $n$ 的表达式为
\begin{equation}
n = \sqrt{NN'}\, e^{-u/2k_{\!B}T},
\end{equation}
其中，$N$ 为晶体中原子的总数，$N'$ 为晶体中的间隙位置总数，$u$ 代表将一个原子从正常格点位移到一个间隙位置而产生的一个弗仑克尔缺陷所需的能量；

(2) 在简单晶格中可以假设间隙位置总数与晶体原子总数（格点总数）近似相等，又一个弗仑克尔缺陷形成能 $u=1\,\mathrm{eV}$，那么要使晶体中有千分之一的原子变成间隙原子需要多高的温度？
\end{modelbox}

\begin{tipbox}[考点快速分析]
\begin{itemize}
    \item \textbf{缺陷类型}：弗仑克尔缺陷 = 空位 + 填隙原子（成对产生），区别于肖特基缺陷（仅空位）
    \item \textbf{统计热力学方法}：用最概然分布（配分函数取极大）或自由能极小求平衡浓度
    \item \textbf{组合数学}：从 $N$ 个格点选 $n$ 个空位，再从 $N'$ 个间隙选 $n$ 个填隙原子
    \item \textbf{斯特林近似}：$\ln x! \approx x\ln x - x$（$x\gg 1$），是熵求导的核心工具
    \item \textbf{低浓度近似}：$n\ll N, n\ll N'$，使 $N-n\approx N$，指数因子中的 $1/2$ 即源于此
\end{itemize}
\end{tipbox}

\begin{derivationbox}[标准解题过程]
\textbf{Step 1: 微观状态数}

弗仑克尔缺陷由一个空位和一个填隙原子组成。设从 $N$ 个正常格点中移走 $n$ 个原子形成 $n$ 个空位，可能的方式数为
\begin{equation}
W_1 = \binom{N}{n} = \frac{N!}{(N-n)!\,n!}.
\end{equation}
这 $n$ 个原子进入 $N'$ 个间隙位置，可能的方式数为
\begin{equation}
W_2 = \binom{N'}{n} = \frac{N'!}{(N'-n)!\,n!}.
\end{equation}
形成 $n$ 个弗仑克尔缺陷的总微观状态数为
\begin{equation}
W(n) = W_1 W_2 = \frac{N!\,N'!}{(N-n)!\,(N'-n)!\,(n!)^2}.
\label{eq:frenkel-W}
\end{equation}

\textbf{Step 2: 能量与配分函数}

产生 $n$ 个缺陷使晶体能量增加 $\Delta U = nu$。缺陷子系统的配分函数为
\begin{equation}
Z = \sum_n W(n)\, e^{-nu/k_{\!B}T}.
\end{equation}
在热平衡时，系统处于最概然分布，对应 $W(n)\,e^{-nu/k_{\!B}T}$ 取极大值。

\textbf{Step 3: 斯特林公式求极值}

对 $W(n)\,e^{-nu/k_{\!B}T}$ 取对数，利用斯特林近似 $\ln x! \approx x\ln x - x$（$x\gg 1$）：
\begin{multline}
\ln\bigl[W(n)\,e^{-nu/k_{\!B}T}\bigr] = N\ln N + N'\ln N' - 2n\ln n \\
- (N-n)\ln(N-n) - (N'-n)\ln(N'-n) - \frac{nu}{k_{\!B}T}.
\end{multline}
对 $n$ 求导，由极值条件 $\partial/\partial n = 0$：
\begin{equation}
-2\ln n + \ln(N-n) + \ln(N'-n) - \frac{u}{k_{\!B}T} = 0.
\end{equation}
整理得
\begin{equation}
\ln\!\left[\frac{n^2}{(N-n)(N'-n)}\right] = -\frac{u}{k_{\!B}T},
\end{equation}
即
\begin{equation}
n = \sqrt{(N-n)(N'-n)}\; e^{-u/2k_{\!B}T}.
\label{eq:frenkel-exact}
\end{equation}

\textbf{Step 4: 低浓度近似}

在通常温度下，缺陷浓度很低，即 $n\ll N$ 且 $n\ll N'$。此时 $N-n\approx N$，$N'-n\approx N'$，由 \eqref{eq:frenkel-exact} 得
\begin{equation}
\boxed{n \approx \sqrt{NN'}\, e^{-u/2k_{\!B}T}.}
\label{eq:frenkel-approx}
\end{equation}
\textit{证毕。}

\textbf{Step 5: 给定缺陷浓度所需温度的计算}

简单晶格中，间隙位置数近似等于格点数，即 $N'=N$。此时弗仑克尔缺陷的平衡浓度为
\begin{equation}
\frac{n}{N} = e^{-u/2k_{\!B}T}.
\end{equation}
题目要求千分之一的原子成为间隙原子，即 $n/N = 10^{-3}$。取对数：
\begin{equation}
\ln\!\left(\frac{n}{N}\right) = -\frac{u}{2k_{\!B}T}
\quad\Longrightarrow\quad
T = \frac{u}{2k_{\!B}\ln(N/n)}.
\end{equation}
代入已知数据：
\begin{align}
u &= 1\,\mathrm{eV} = 1.602\times 10^{-19}\,\mathrm{J}, \\
k_{\!B} &= 1.381\times 10^{-23}\,\mathrm{J/K}, \\
\ln(N/n) &= \ln 1000 \approx 6.9078,
\end{align}
得
\begin{equation}
T = \frac{1.602\times 10^{-19}}{2\times 1.381\times 10^{-23}\times 6.9078}
= \frac{1.602\times 10^{-19}}{1.907\times 10^{-22}}
\approx \boxed{840\,\mathrm{K}}.
\end{equation}
\end{derivationbox}

\begin{formulabox}[答案汇总]
\begin{center}
\begin{tabular}{ll}
\toprule
\textbf{问题} & \textbf{答案} \\
\midrule
(1) 弗仑克尔缺陷数目 & $n = \sqrt{NN'}\, e^{-u/2k_{\!B}T}$ \\
(2) 所需温度 & $T \approx 840\,\mathrm{K}$ \\
\bottomrule
\end{tabular}
\end{center}
\end{formulabox}

\begin{insightbox}[物理图像：指数上的 $1/2$ 从何而来？]
弗仑克尔缺陷由空位和填隙原子\textbf{成对}产生，涉及两个独立的组合过程：
\begin{itemize}
    \item 从 $N$ 个格点中选 $n$ 个空位：熵贡献 $S_1 = k_{\!B}\ln W_1$
    \item 从 $N'$ 个间隙中选 $n$ 个填隙原子：熵贡献 $S_2 = k_{\!B}\ln W_2$
\end{itemize}
对自由能 $F = U - TS$ 求极小，两个配置熵的导数各贡献一个 $\ln n$ 项，合计 $-2\ln n$。能量项只贡献一次 $nu$，故平衡时
\[
2k_{\!B}T\ln n \sim u \quad\Longrightarrow\quad n\sim e^{-u/2k_{\!B}T}.
\]
\textbf{对比}：肖特基缺陷（仅空位，无填隙原子）只有一个组合过程，故 $n = N e^{-u/k_{\!B}T}$，指数上没有 $1/2$。
\end{insightbox}

\begin{thinkbox}[考场速记：缺陷统计问题的答题模板]
遇到``缺陷平衡浓度''类问题，按以下步骤展开：
\begin{enumerate}
    \item \textbf{写微观状态数}：区分缺陷类型（弗仑克尔/肖特基/反位），写出组合数 $W(n)$
    \item \textbf{写能量和自由能}：$F(n) = nu - k_{\!B}T\ln W(n)$
    \item \textbf{求极值}：$\partial F/\partial n = 0$，或用 $\ln W(n) - nu/k_{\!B}T$ 取极大
    \item \textbf{用斯特林近似}：$\partial(x\ln x)/\partial x = \ln x + 1$，常数项在求导时消去
    \item \textbf{低浓度近似}：$n\ll N$ 时 $N-n\approx N$，得到最终简洁形式
    \item \textbf{检验指数}：
    \begin{itemize}
        \item 弗仑克尔（成对）：$e^{-u/2k_{\!B}T}$
        \item 肖特基（单空位）：$e^{-u/k_{\!B}T}$
    \end{itemize}
\end{enumerate}
\textit{陷阱提醒}：不要混淆 $N$（格点数）和 $N'$（间隙位置数）。简单晶格中可假设 $N'=N$，但体心立方/面心立方的间隙位置数不同（八面体间隙、四面体间隙数目不同），此时 $N'\neq N$。
\end{thinkbox}
